% Compile with XeLaTeX. This file has no external assets or dependencies beyond TeX packages.
\documentclass[11pt,a4paper]{article}
\usepackage{fontspec}
\usepackage{polyglossia}
\setmainlanguage{ukrainian}
\setmainfont{DejaVu Serif}
\setsansfont{DejaVu Sans}
\setmonofont{DejaVu Sans Mono}
\usepackage[a4paper,margin=22mm]{geometry}
\usepackage{amsmath,amssymb,mathtools,booktabs,array}
\usepackage{microtype}
\usepackage{xurl}
\setlength{\emergencystretch}{3em}
\usepackage[unicode,colorlinks=true,linkcolor=blue,urlcolor=blue]{hyperref}
\usepackage{fancyhdr}
\pagestyle{fancy}\fancyhf{}
\fancyhead[L]{\small Метрологія · Лекція 9}
\fancyhead[R]{\small Математичний конспект}
\fancyfoot[C]{\thepage}
\setlength{\headheight}{14pt}
\setlength{\parindent}{0pt}
\setlength{\parskip}{5pt}
\newcommand{\K}{\mathrm{K}}
\newcommand{\C}{\,{}^{\circ}\mathrm{C}}
\newcommand{\Ohm}{\,\Omega}
\newcommand{\unit}[1]{\,\mathrm{#1}}
\newcommand{\defn}[1]{\paragraph{#1}}
\title{Температурні вимірювання\newline та вибір термометра\\[6pt]\large Лекція 9. Математичний конспект}
\author{Метрологія, технологічні вимірювання та прилади}
\date{}
\begin{document}
\raggedright
\maketitle
\vspace{-4mm}
\textbf{Позначення.} $T$ --- термодинамічна температура в кельвінах;
$t$ --- температура Цельсія; $s$ --- час у секундах.
Індекс $\infty$ позначає усталене значення, індекс $0$ --- опорне або початкове значення за явно заданих умов.
Відношення $T/\K$ позначає числове значення температури в кельвінах.
Температурний інтервал $1\C$ за розміром дорівнює $1\K$.

\section{Температурні шкали й термометри розширення}

\defn{Термодинамічна температура та шкала Цельсія}
Одиницю кельвін визначено фіксацією сталої Больцмана
\[
 k_{\mathrm B}=1.380649\cdot10^{-23}\unit{J\,K^{-1}}
 \quad\text{(точно)}.
\]
Зв'язок числових значень шкал і його обернення:
\begin{align}
 \frac{T}{\K}&=\frac{t}{{}^{\circ}\mathrm C}+273.15, &
 \frac{t}{{}^{\circ}\mathrm C}&=\frac{T}{\K}-273.15,\\
 \frac{\Delta T}{\K}&=\frac{\Delta t}{{}^{\circ}\mathrm C}.
\end{align}
Для звичайних теплових станів у цій лекції $T>0$. Температуру $T_{90}$, визначену за ITS-90, відрізняють від $T$: практична шкала реалізує визначені процедури й інтерполяцію, а різниця $T-T_{90}$ не є тотожно нульовою.

\defn{Лінійне розширення}
Для вільного однорідного стрижня за малих деформацій та майже сталого коефіцієнта $\alpha_\ell$:
\begin{align}
 \ell(T)&\approx\ell_0[1+\alpha_\ell(T-T_0)], &
 \Delta\ell&\approx\alpha_\ell\ell_0\Delta T,\\
 \Delta T&\approx\frac{\Delta\ell}{\alpha_\ell\ell_0}, &
 [\alpha_\ell]&=\K^{-1},\quad [\ell_0]=[\Delta\ell]=\mathrm m.
\end{align}
Обернення потребує $\alpha_\ell\ell_0\ne0$.
Умова малості: $|\alpha_\ell\Delta T|\ll1$.
Якщо $\alpha_\ell$ істотно залежить від температури, локальне означення має вигляд
\[
 \alpha_\ell(T)=\frac{1}{\ell(T)}\frac{d\ell}{dT},\qquad
 \ln\frac{\ell(T)}{\ell(T_0)}=\int_{T_0}^{T}\alpha_\ell(u)\,du.
\]
Для рідини вводять коефіцієнт об'ємного розширення $\beta$ з одиницею $\K^{-1}$:
\[
 \Delta V\approx\beta V_0\Delta T.
\]
У спрощеній капілярній моделі з площею перерізу $A_c$:
\[
 \Delta h\approx\frac{V_0}{A_c}(\beta_{\rm liquid}-\beta_{\rm vessel})\Delta T,
 \quad [V_0]=\mathrm{m^3},\ [A_c]=\mathrm{m^2},\ [\Delta h]=\mathrm m.
\]
Модель враховує різницю розширень заповнення й оболонки, але нехтує зміною площі капіляра та неоднорідністю температури.

\defn{Біметалевий і газовий манометричний принципи}
Для двох незв'язаних шарів однакової початкової довжини різниця вільних видовжень
\[
 \Delta\ell_1-\Delta\ell_2\approx
 (\alpha_1-\alpha_2)\ell_0\Delta T.
\]
У зв'язаній пластині ця несумісність деформацій спричиняє кривизну, що залежить також від товщин і модулів пружності; наведений вираз не є формулою кута повороту.

Для ідеального газу сталої кількості $n$ у жорсткому об'ємі $V$:
\[
 pV=nR_uT,\qquad \frac{p}{p_0}=\frac{T}{T_0},\qquad
 T=T_0\frac{p}{p_0}.
\]
$p,p_0$ --- \emph{абсолютні} тиски в паскалях, $V$ --- м$^3$,
$n$ --- моль, $R_u$ --- Дж/(моль$\cdot$К).
Модель вимагає однорідної температури заповнення, сталого об'єму й кількості газу та не описує довільний рідинний або парорідинний манометричний термометр.

\defn{Тепловий баланс контактного сенсора}
Модель зосередженої теплоємності:
\begin{equation}
 C_{\rm th}\frac{dT_s}{ds}
 =G_o(T_o-T_s)+G_a(T_a-T_s)+P.
\end{equation}
Тут $C_{\rm th}$ --- Дж/К; $G_o,G_a$ --- теплові провідності до об'єкта й оточення, Вт/К; $P$ --- внутрішнє тепловиділення, Вт; усі $T$ --- К; $s$ --- с.
За сталих коефіцієнтів і температур:
\begin{align}
 T_{s,\infty}&=\frac{G_oT_o+G_aT_a+P}{G_o+G_a}, &
 \tau&=\frac{C_{\rm th}}{G_o+G_a},\\
 T_{s,\infty}-T_o&=\frac{G_a(T_a-T_o)+P}{G_o+G_a}.&&
\end{align}
Тому усталення показу не доводить рівності температур сенсора та об'єкта.

\defn{Перехідний процес і критерій усталення}
За одного сталого кінцевого значення $T_\infty$, $\tau>0$ та початкової температури $T_0$:
\begin{align}
 \tau\frac{dT_s}{ds}+T_s&=T_\infty, &T_s(0)&=T_0,\\
 T_s(s)&=T_\infty+(T_0-T_\infty)e^{-s/\tau}, & s&\ge0.
\end{align}
Одиниці: $T_s,T_0,T_\infty$ --- К; $s,\tau$ --- с.
Для абсолютної смуги $\varepsilon_T>0$ (К) час усталення визначається як
\begin{equation}
 s_{\rm set}=\inf\left\{s_0\ge0:
 |T_s(s)-T_\infty|\le\varepsilon_T
 \quad\text{для всіх }s\ge s_0\right\}.
\end{equation}
Інфімум порожньої множини дорівнює $+\infty$. Для експоненційної моделі
\[
 s_{\rm set}=\begin{cases}
 0,& |T_0-T_\infty|\le\varepsilon_T,\\
 \tau\ln\dfrac{|T_0-T_\infty|}{\varepsilon_T},& |T_0-T_\infty|>\varepsilon_T.
 \end{cases}
\]
Для ненульового стрибка й відносної смуги $0<\eta<1$,
$\varepsilon_T=\eta|T_0-T_\infty|$, тому
\[
 s_{\rm set}=-\tau\ln\eta.
\]
За $\eta=0.02$: $s_{\rm set}=\tau\ln50\approx3.912023\tau$.
Після $3\tau$ залишок становить $e^{-3}\approx0.049787$ початкової різниці.
Для коливального відгуку перше входження у смугу недостатнє.
Скінченний дискретний запис без додаткової моделі перевіряє смугу лише у спостережених точках, не в усіх майбутніх моментах.

\section{Терморезистивні перетворювачі}

\defn{Металевий термометр опору}
Локальна лінійна модель за опорної температури $t_0$:
\begin{align}
 R(t)&\approx R_0[1+\alpha_R(t-t_0)],\\
 t-t_0&\approx\frac{R/R_0-1}{\alpha_R}.
\end{align}
$R,R_0$ --- Ом; $\alpha_R$ --- К$^{-1}$; $t-t_0$ --- температурний інтервал у К (чисельно такий самий у $^{\circ}$C).
Обчислення числового значення $t$ в $^{\circ}$C:
\[
 \frac{t}{{}^{\circ}\mathrm C}
 =\frac{t_0}{{}^{\circ}\mathrm C}
 +\frac{R/R_0-1}{\alpha_R/(\K^{-1})}.
\]
Потрібні $R_0>0$, $\alpha_R\ne0$ та $t$ у прийнятому інтервалі моделі.
Чутливість $S_R=dR/dt\approx R_0\alpha_R$ має одиницю Ом/К.
У навчальній моделі $R_0=100\Ohm$, $t_0=0\C$,
$\alpha_R=0.00385\unit{K^{-1}}$, $0\C\le t\le100\C$:
\[
 R(0\C)=100\Ohm,\qquad R(100\C)=138.5\Ohm.
\]
Ці точки належать лінійному наближенню, а не точній універсальній характеристиці реального Pt100.
Номінальне позначення Pt100 задає $100\Ohm$ за $0\C$, але не клас допуску, струм або конструктивний діапазон.

\defn{Нелінійна характеристика платини}
Поширена форма рівняння Каллендара--Ван Дюзена:
\[
 R(t)=\begin{cases}
 R_0(1+At+Bt^2),& t\ge0\C,\\
 R_0[1+At+Bt^2+C(t-100\C)t^3],&t<0\C.
 \end{cases}
\]
Тут $t$ записується як температура Цельсія відносно $0\C$,
а одиниці коефіцієнтів узгоджені з цим записом:
$A$ --- $({}^{\circ}\mathrm C)^{-1}$,
$B$ --- $({}^{\circ}\mathrm C)^{-2}$,
$C$ --- $({}^{\circ}\mathrm C)^{-4}$.
Конкретні коефіцієнти та допустимий інтервал задають для відповідної характеристики або за калібруванням; це не продовження навчальної прямої без обмежень.

\defn{Двопровідне вимірювання}
Для послідовних опорів проводів $r_1,r_2$:
\begin{align}
 R_{\rm meas}&=R_s+r_1+r_2, &
 R_s&=R_{\rm meas}-r_1-r_2,\\
 t_{\rm raw}-t_{\rm corr}&=\frac{r_1+r_2}{R_0\alpha_R}.
\end{align}
Всі опори --- Ом, температурна різниця --- К.
Рівність останнього рядка точна всередині прийнятої лінійної моделі.
За $r_1,r_2\ge0$, $\alpha_R>0$ зміщення невід'ємне.
За нульового опору проводів обидві температури збігаються.
Корекція вважає $r_1,r_2$ відомими при фактичних умовах вимірювання.

\defn{Міст і розділення струмових та потенціальних кіл}
Для двох ідеальних резистивних подільників з верхніми опорами $R_1,R_3$,
нижніми $R_2,R_4$ та спільним живленням $U_{\rm exc}$:
\[
 U_o=U_{\rm exc}\left(\frac{R_2}{R_1+R_2}-\frac{R_4}{R_3+R_4}\right),
 \qquad U_o=0\iff R_2R_3=R_4R_1.
\]
Напруги --- В, опори --- Ом; вимірювач виходу вважається ідеально високоомним.
Трипровідна компенсація залежить від конкретної мостової або вимірювальної схеми й узгодженості відповідних опорів проводів.
У чотирипровідній схемі за відомого струму $I$ та нехтовно малого струму потенціальних входів:
\[
 U_s=IR_s,\qquad R_s=\frac{U_s}{I},\qquad [U_s]=\mathrm V,\ [I]=\mathrm A.
\]
Скінченні вхідні струми, паразитні ЕРС та невизначеність $I$ залишаються джерелами відхилень.

\defn{Самонагрівання}
\begin{align}
 P&=I^2R_s, & \Delta T_{\rm self}&\approx\theta P,\\
 [P]&=\mathrm W, & [\theta]&=\K/\mathrm W.
\end{align}
Умова: усталений режим, майже сталі $R_s$ і тепловий опір $\theta$.
Якщо задано коефіцієнт розсіювання $D$ у Вт/К, то $\theta=1/D$ та
$\Delta T_{\rm self}=P/D$. Одиниці Вт/К і К/Вт не можна переставляти.
За незмінних $R_s,\theta$ множення струму на $a$ множить нагрівання на $a^2$.
Корекція опору проводів цього нагрівання не усуває.

\defn{NTC-термістор: B-модель}
\begin{align}
 R(T)&=R_0\exp\left[B_N\left(\frac1T-\frac1{T_0}\right)\right],\\
 \frac1T&=\frac1{T_0}+\frac1{B_N}\ln\frac{R}{R_0},\\
 T&=\left[\frac1{T_0}+\frac1{B_N}\ln\frac{R}{R_0}\right]^{-1}.
\end{align}
$T,T_0,B_N$ --- К; $R,R_0$ --- Ом; $T,T_0,R,R_0,B_N>0$.
Для обернення вираз у квадратних дужках має бути додатним.
Показник експоненти та аргумент логарифма безрозмірні.
\[
 \frac{dR}{dT}=-\frac{B_N}{T^2}R<0.
\]
B-модель є наближенням у заданому інтервалі, а не універсальним описом усіх термісторів; PTC-елементи не підпорядковуються наведеній спадній характеристиці.

\section{Термоелектричні перетворювачі}

\defn{Термоелектрична характеристика}
Нехай $S_{AB}(T)$ --- різниця коефіцієнтів Зеєбека двох однорідних провідників, В/К.
Для вибраної полярності:
\[
 E(T_h,T_r)=\int_{T_r}^{T_h}S_{AB}(u)\,du,
 \qquad [E]=\mathrm V,\quad [T_h]=[T_r]=\K.
\]
Якщо $S_{AB}$ приблизно сталий лише на заданому інтервалі, можна використовувати $E\approx S_{AB}(T_h-T_r)$.
Це наближення не замінює таблиці конкретного типу термопари на всьому діапазоні.

\defn{Компенсація опорного з'єднання}
Введемо функцію $E_0(t)=E(t,0\C)$ для відповідного типу й полярності.
\begin{align}
 E_{\rm meas}&=E_0(t_h)-E_0(t_r),\\
 E_{\rm comp}&=E_{\rm meas}+E_0(t_r),\\
 t_h&=E_0^{-1}(E_{\rm comp}).
\end{align}
Температури $t_h,t_r$ --- $^{\circ}$C; усі напруги повинні мати однакову одиницю, наприклад мВ.
Обернена функція використовується лише на однозначній гілці характеристики.
Перевірки:
\[
 E_0(0\C)=0,\qquad
 E_{\rm meas}=0\Longrightarrow t_h=t_r
 \quad\text{на обраній однозначній гілці}.
\]
Для нелінійної $E_0$ загалом
\[
 E_0^{-1}\bigl(E_{\rm meas}+E_0(t_r)\bigr)
 \ne E_0^{-1}(E_{\rm meas})+t_r.
\]
Рівність додавання температур можлива для лінійної характеристики $E_0(t)=at$ із відповідним нулем відліку.
Умови: правильна полярність, однорідні провідники, узгоджені подовжувальні матеріали та враховані температури переходів до вимірювального кола.

\defn{Кусочно-лінійна інтерполяція}
Нехай на одному інтервалі $x_i<x_{i+1}$ задано $f_i,f_{i+1}$.
Для $x\in[x_i,x_{i+1}]$:
\begin{align}
 \lambda&=\frac{x-x_i}{x_{i+1}-x_i}, &0&\le\lambda\le1,\\
 L_i(x)&=(1-\lambda)f_i+\lambda f_{i+1}.
\end{align}
Якщо $x$ --- температура в $^{\circ}$C, то $f$ --- напруга в мВ, а нахил має одиницю мВ/К.
Для строго монотонної гілки таблиці з $E_{i+1}\ne E_i$:
\[
 t(E)=t_i+\frac{E-E_i}{E_{i+1}-E_i}(t_{i+1}-t_i).
\]
Частка напруг безрозмірна; результат --- $^{\circ}$C.
Граничні значення: $t(E_i)=t_i$, $t(E_{i+1})=t_{i+1}$.
Пряме й обернене перетворення на одному відрізку відновлюють аргумент у точній арифметиці.
За межами таблиці інтерполяція не визначена; її продовження є окремою процедурою екстраполяції.

\defn{Оцінка похибки інтерполяції та пошуку}
Якщо $f\in C^2[x_i,x_{i+1}]$ та $|f''(x)|\le M_2$, то
\[
 |f(x)-L_i(x)|\le\frac{M_2}{8}(x_{i+1}-x_i)^2.
\]
Якщо $x$ у $^{\circ}$C і $f$ у мВ, $M_2$ має одиницю мВ/К$^2$.
Для оберненої функції, якщо $|E_0'(t)|\ge m>0$ (мВ/К), локальне відхилення напруги $\delta E$ обмежує відхилення температури:
\[
 |\delta t|\le\frac{|\delta E|}{m}
\]
на інтервалі, де обидва аргументи належать одній гілці та виконується нижня межа похідної.
Ця оцінка не є повним бюджетом невизначеності термопарного каналу.
Послідовний пошук серед $n$ вузлів має часову складність $O(n)$; після знаходження інтервалу інтерполяція потребує $O(1)$ арифметичних операцій.

\section{Випромінювання, акустика і вибір методу}

\defn{Сумарне випромінювання}
Для чорного тіла густина власного потоку випромінювання
$M_b=\sigma T^4$. Для сірої поверхні:
\begin{equation}
 M=\varepsilon\sigma T^4,
 \quad\sigma\approx5.670374419\cdot10^{-8}\unit{W\,m^{-2}\,K^{-4}}.
\end{equation}
$M$ --- Вт/м$^2$; $T$ --- К; $0<\varepsilon\le1$ --- безрозмірна випромінювальна здатність, прийнята незалежною від довжини хвилі.
Якщо відомий саме власний сумарний потік:
\[
 T=\left(\frac{M}{\varepsilon\sigma}\right)^{1/4}.
\]
Це не є універсальним перетворенням показу спектрального пірометра.

\defn{Власний та відбитий потоки}
Для непрозорої сірої дифузної поверхні в ізотропному ізотермічному чорному оточенні:
\begin{align}
 J&=\varepsilon\sigma T^4+(1-\varepsilon)\sigma T_b^4,\\
 T&=\left[\frac{J-(1-\varepsilon)\sigma T_b^4}{\varepsilon\sigma}\right]^{1/4}.
\end{align}
$J$ --- сумарна вихідна густина потоку, Вт/м$^2$; $T_b$ --- температура фону, К.
Для обернення потрібні $\varepsilon>0$ та додатний підкореневий вираз.
За $T=T_b$ отримуємо $J=\sigma T^4$, незалежно від $\varepsilon$ у цій моделі.
За $\varepsilon\to0$ залежність від температури самої поверхні зникає, тому обернення стає нестійким.

\defn{Спектральна модель приймача}
Спектральна енергетична яскравість чорного тіла:
\[
 L_\lambda^b(T)=\frac{2hc_0^2}{\lambda^5}
 \left[\exp\left(\frac{hc_0}{\lambda k_{\rm B}T}\right)-1\right]^{-1}.
\]
$\lambda$ --- м; $c_0$ --- швидкість світла, м/с; $h$ --- Дж$\cdot$с;
$L_\lambda^b$ --- Вт/(м$^2\cdot$ср$\cdot$м).
Спрощена модель сигналу для непрозорої поверхні, повного заповнення плями й без атмосферного поглинання:
\[
 S=C_g\int_{\lambda_1}^{\lambda_2}w(\lambda)
 \left[\varepsilon_\lambda L_\lambda^b(T)
 +(1-\varepsilon_\lambda)L_\lambda^b(T_b)\right]d\lambda.
\]
$w$ --- безрозмірна спектральна вагова функція;
$C_g$ містить геометрію та коефіцієнт перетворення приймача;
для $S$ у вольтах $[C_g]=\mathrm{V\,m^2\,sr/W}$.
Неповне заповнення плями, пропускання оптики, поглинання та власне випромінювання шляху потребують додаткових членів.

\defn{Швидкість звуку в ідеальному газі}
Для однорідного газу сталого складу за малих адіабатичних збурень:
\begin{align}
 c_a^2&=\gamma\frac{R_u}{M_{\rm mol}}T, &
 T&=\frac{M_{\rm mol}c_a^2}{\gamma R_u},\\
 \gamma&=\frac{c_p}{c_v}>1, &
 R_u&\approx8.314462618\unit{J\,mol^{-1}\,K^{-1}}.
\end{align}
$c_a$ --- м/с; $M_{\rm mol}$ --- кг/моль; $\gamma$ --- безрозмірне;
$T$ --- К. Розмірнісна перевірка:
\[
 \frac{(\mathrm{kg/mol})(\mathrm{m^2/s^2})}
 {\mathrm{J/(mol\,K)}}=\K.
\]
Сталу $\gamma$ можна застосовувати лише в прийнятому температурному інтервалі.
Залежність складу, неідеальність газу та дисперсія можуть вимагати інших моделей.

\defn{Час прольоту та обернення}
Для нерухомого однорідного газу та \emph{однобічної} бази $L$:
\begin{align}
 c_a&=\frac{L}{\Delta s}, &
 T&=\frac{M_{\rm mol}L^2}{\gamma R_u(\Delta s)^2},\\
 \Delta s&=L\sqrt{\frac{M_{\rm mol}}{\gamma R_uT}}.
\end{align}
$L>0$ --- м; $\Delta s>0$ --- с; $M_{\rm mol}>0$ --- кг/моль.
Якщо задано час у мілісекундах, $\Delta s/\mathrm s=10^{-3}(\Delta s/\mathrm{ms})$.
Для ехо-сигналу з відстанню до відбивача $L$:
\[
 c_a=\frac{2L}{\Delta s_{\rm echo}}.
\]
Використання $L$ замість $2L$ зменшує швидкість удвічі та оцінку абсолютної температури вчетверо.

\defn{Обмеження діапазону та чутливість}
Для допустимого інтервалу $0<T_{\min}\le T\le T_{\max}$:
\[
 L\sqrt{\frac{M_{\rm mol}}{\gamma R_uT_{\max}}}
 \le\Delta s\le
 L\sqrt{\frac{M_{\rm mol}}{\gamma R_uT_{\min}}}.
\]
Границі включено. За сталих інших параметрів
\[
 \frac{T'}T=\left(\frac{L'}L\right)^2,\qquad
 \frac{T'}T=\left(\frac{\Delta s}{\Delta s'}\right)^2,\qquad
 \frac{\partial T}{\partial(\Delta s)}=-\frac{2T}{\Delta s}.
\]
Остання похідна має одиницю К/с. Ці співвідношення оцінюють вплив одного параметра за фіксованих інших, а не повний бюджет невизначеності.

\defn{Потік і неоднорідність}
Для однорідного газу зі швидкістю $v$ уздовж бази, $|v|<c_a$:
\[
 \Delta s_+=\frac{L}{c_a+v},\qquad
 \Delta s_-=\frac{L}{c_a-v},\qquad
 c_a=\frac L2\left(\frac1{\Delta s_+}+\frac1{\Delta s_-}\right).
\]
Для нерухомого газу з неоднорідною температурою вздовж прямого шляху:
\[
 \Delta s=\int_0^L\frac{dx}{c_a(T(x))}.
\]
Один час не визначає довільний профіль $T(x)$ і загалом не дає його арифметичного середнього.
Це обмеження важливе при зіставленні акустичного результату з локальним контактним вимірюванням.

\defn{Формальна придатність каналу}
Для заданої вимірюваної величини придатність вимагає одночасного виконання
\[
 D_{\rm required}\subseteq D_{\rm model}\cap D_{\rm instrument},
 \qquad \text{умов монтажу, середовища й динаміки}.
\]
Порушення відомої обов'язкової умови означає непридатність.
Відсутність потрібного параметра за відсутності встановленої відмови означає недостатність даних, а не підтверджену придатність.

\paragraph{Нормативні й технічні опори формул}
\small
\begin{itemize}
\item BIPM, визначення кельвіна та практична реалізація:
\url{https://www.bipm.org/en/si-base-units/kelvin}.
\item JCGM VIM 3, 4.23 (час відгуку на стрибок):
\url{https://jcgm.bipm.org/vim/en/4.23.html}.
\item Analog Devices, RTD Measurement System Design Essentials:
\url{https://www.analog.com/en/resources/technical-articles/rtd-measurement-system-design-essentials.html}.
\item Analog Devices, CN0394 (компенсація опорного з'єднання):
\url{https://www.analog.com/en/resources/reference-designs/circuits-from-the-lab/CN0394.html}.
\item NIST, Acoustic Thermometry:
\url{https://www.nist.gov/programs-projects/acoustic-thermometry}.
\end{itemize}
\end{document}
